\chapter{Processus métier bout-en-bout}

\section{P1 --- Inscription entreprise et activation}

\subsection{Étapes produit}

\begin{enumerate}
  \item Création compte IAM (signup BFF $\rightarrow$ Core $\rightarrow$ Kernel).
  \item Formulaire multi-étapes \pathref{app/register/page.tsx}
        (identité, KYC docs, settings, credentials).
  \item \texttt{registerService.submitApplication} $\rightarrow$
        \api{/onboarding/applications}.
  \item Complétion identité Kernel :
        \texttt{completeKernelIdentity} / \texttt{ensureKernelIdentity}.
  \item Attente \api{/pending} (\texttt{getMyApplication}).
  \item Admin TNT : \api{/admin/onboarding/[id]} --- approve/reject
        (\pathref{onboardingAdminService.ts}).
\end{enumerate}

\subsection{Contrainte Kernel critique}

\pathref{ONBOARDING\_KERNEL\_FLOW.md} : l'endpoint actors/onboarding est lié au
\texttt{sub} du JWT. La phase 1 \textbf{doit} s'exécuter avec le JWT
\textbf{candidat}. Approuver avec le JWT admin sans phase 1 associée crée le
mauvais BusinessActor.

\subsection{Phase admin (orchestration Core)}

Approve $\rightarrow$ organisation Kernel + rôles TNT + options plateforme
(trust ON) + rôle manager. Prérequis : \texttt{kernelIdentityReady=true}.

SSO post-activation (HRM / KSM\ldots) exige \texttt{kernelOrganizationId}.

\begin{lstlisting}[caption={Séquence onboarding (simplifiée)}]
Candidat -> BFF signup -> Core auth/sign-up -> Kernel
Candidat -> BFF onboarding/applications (+ Bearer candidat)
Candidat -> BFF .../kernel-identity
Admin    -> BFF admin/.../approve -> Core onboarding approve
         -> Agency ACTIVE + org Kernel liée
\end{lstlisting}

\section{P2 --- Intake client et création de mission}

\subsection{Canal QR public}

\begin{enumerate}
  \item HQ génère URL : \texttt{intakeDepositUrl(agencyId, branchId)}
        $\rightarrow$ \api{/track/deposit?agencyId=\&branchId=}.
  \item Client charge contexte : \api{GET agencies/\{id\}/intake-context}.
  \item Submit : \api{POST intake-requests}.
  \item Agent HQ : file \api{/accueil/demandes} --- approve/reject.
  \item Walk-in direct : \api{/accueil} via \texttt{intakeService.walkIn}.
\end{enumerate}

\subsection{Canal suivi}

\api{/track} : recherche code, stream SSE optionnel, claims / drop-off / hubs
relais (\pathref{trackingService.ts}, \pathref{complianceService.ts}).

\section{P3 --- Lifecycle mission}

Statuts (\pathref{lib/types.ts}) :
\texttt{DRAFT → PENDING → ASSIGNED → IN\_TRANSIT → AT\_HUB → DELIVERED}
(ou \texttt{FAILED}/\texttt{CANCELLED}).

\begin{description}[style=nextline]
  \item[HQ / Antenne]
    create, assign/reassign, cancel/reschedule, deposit-hub, withdraw
    (\pathref{missionService.ts}, services branch).
  \item[Livreur]
    pickup, deliver (+ POD), anomaly ; si offline $\rightarrow$ queue IndexedDB
    puis sync (\pathref{livreurMissionService}, \pathref{offlineQueue.ts}).
\end{description}

\section{P4 --- Facturation et commissions}

\begin{itemize}
  \item Politiques tarifaires actives (devise XAF typique).
  \item Estimation à la création de mission.
  \item Factures / revenus après clôture.
  \item Commissions livreur (portail gains + HQ personnel).
\end{itemize}
Services : \pathref{billingService.ts}, \pathref{livreurCommissionService.ts}.

\section{P5 --- Flotte et FleetMan}

\begin{enumerate}
  \item CRUD véhicules Agency (registry).
  \item Pont FleetMan (managers HQ) :
        \api{agencies/\{id\}/fleetman/\{status|connect|launch|sync\}}.
  \item Tokens stockés chiffrés AES
        (\envvar{TNT\_FLEETMAN\_TOKEN\_ENCRYPTION\_KEY}) ;
        plaintext interdit en prod.
\end{enumerate}

\section{P6 --- Hubs et Trust}

Transitions hub / POD peuvent ancrer des preuves côté Trust \textbf{via Core}
après activation. Le frontend expose
\envvar{NEXT\_PUBLIC\_TNT\_TRUST\_BASE\_URL} mais n'implémente pas (audit) un
client Trust autonome --- la vérité passe par les APIs agency/Core.
